\chapter{Introducere}
\label{cap:introducere}

Pasiunea pentru domeniul modei și estetica vizuală a reprezentat întotdeauna un punct de interes personal. Observația constantă a modului în care oamenii se raportează la vestimentație a scos la iveală o provocare comună: dificultatea de a compune ținute coerente și de a se alinia la tendințele actuale. Multe persoane, deși dețin o garderobă variată, se confruntă zilnic cu indecizia și nesiguranța, neavând la dispoziție cunoștințele unui designer vestimentar.

Această lucrare de licență pornește de la această problemă concretă. Proiectul dezvoltă o soluție tehnologică, sub forma unei aplicații mobile, care funcționează ca un asistent stilistic personal. Aplicația oferă utilizatorilor recomandări de ținute bazate pe articolele vestimentare proprii, sugerând combinații de culori armonioase și haine care corespund tendințelor actuale. Astfel, procesul de a te îmbrăca devine mai simplu, mai rapid și mai creativ.

\section{Contextul Tehnic și Științific}

Industria modei și modul în care utilizatorii interacționează cu produsele vestimentare au fost puternic influențate de digitalizare și de migrarea serviciilor către platforme online și aplicații mobile. În prezent, utilizatorii se confruntă cu un volum mare de opțiuni (atât în mediul de e-commerce, cât și în propria garderobă), iar procesul de alegere a unei ținute potrivite devine adesea dificil: necesită timp, presupune decizii repetate și implică atât preferințe personale, cât și criterii de compatibilitate vizuală (culori, contrast, echilibru, stil) \cite{kang2017visually}. Această problemă, cunoscută în literatură sub termenul de "decision fatigue", reduce satisfacția utilizatorului și impactează negativ experiența de e-commerce \cite{fashion_ai_review}.

În acest context, sistemele de recomandare au devenit o componentă esențială pentru personalizarea experienței utilizatorului, fiind utilizate pentru a sugera elemente relevante în funcție de comportament, preferințe și context. Putem identifica mai multe direcții majore pentru construirea recomandărilor: metode \textit{content-based}, metode \textit{collaborative} și abordări hibride, fiecare având avantaje și limitări specifice \cite{adomavicius2005toward}. În practică, în special pentru domenii cu multe opțiuni și gusturi subiective (precum moda), personalizarea este un factor important pentru creșterea satisfacției utilizatorului și a calității experienței \cite{su2009survey}.

Un subset important al domeniului îl reprezintă recomandarea sau generarea unei ținute complete (outfit), unde obiectivul nu este doar selectarea unui singur produs, ci alegerea unei combinații coerente de articole \cite{mcauley2015image}. Alegerea unei combinații poate fi abordată fie prin modele de preferință învățate din date (abordare machine learning), fie prin euristici și reguli explicabile, construite pe baza unor principii de compatibilitate (de exemplu compatibilitatea cromatică) \cite{han2017learning}. În plus, în modă, relevanța sugestiilor este influențată și de tendințe sezoniere/anuale (culori populare, stiluri), astfel că introducerea unui „semnal de trend" poate crește utilitatea practică a recomandărilor \cite{kang2017visually}.

\section{Motivația și Relevanța Problemei}

Domeniul fashion technology (Fashion AI) a cunoscut o evoluție accelerată în ultimii 5 ani, cu aplicații precum Acloset, Whering și Polyvore dominând piața retail vestimentar \cite{fashion_ai_review}. Cu toate acestea, aceste soluții prezintă limitări în ceea ce privește transparența algoritmului, flexibilitatea regulilor de compatibilitate și integrarea efectivă a semnalelor de trend.

Relevanța acestei lucrări se manifestă pe mai multe dimensiuni:

\begin{enumerate}
    \item \textbf{Transparență Algoritmică}: Spre deosebire de metodele black-box bazate pe deep learning, acest proiect propune o abordare explicabilă unde utilizatorul înțelege rațiunea din spatele fiecărei recomandări.
    \item \textbf{Scalabilitate Mobilă}: Aplicația este optimizată pentru dispozitivele mobile cu constrângeri de resurse, utilizând React Native pentru maximă compatibilitate cross-platform \cite{reactnative_doc}.
    \item \textbf{Integrare Trend Semantică}: În contrast cu competitorii care ignorează sau hardcodează trendurile, StyleGenAI oferă o metodologie pentru integrare dinamică a semnalelor sezoni